\chapter{核心数据流}

\section{认知形成}

\begin{enumerate}
\item 宿主或外部系统保存原始材料或规范外部事实。
\item 主体实际接触相应内容时形成观察事件。
\item 证据引用固定来源区域、派生表示和来源谱系。
\item 确定性流程或模型产生认知提案，并声明目标领域、主张、认知指向和输入修订。
\item 目标领域所有者核验证据、权限、对象状态和并发版本。
\item 通过核验后提交不可变修订，推进对象头和权威提交序号。
\item 权威变更使相关服务投影失效或进入增量更新流程。
\end{enumerate}

\section{认知检索}

\begin{enumerate}
\item 调用方提交认知查询和计算预算。
\item 系统解析硬约束、软偏好、主体范围和调用方能力。
\item 查询计划在候选生成前固定检索通道、预算和融合参数。
\item 各检索通道生成服务视图命中，并在对象修订层聚合。
\item 跨通道融合和可选重排产生候选顺序。
\item 返回前重新核验权威修订、生命周期、访问策略和来源状态。
\item 实际进入调用方上下文的对象可以记录呈现事件；进一步引用、行动或结果另行记录类型化使用事件。
\end{enumerate}

\section{离线认知}

\begin{enumerate}
\item 宿主授予维护租约。
\item Nous Wave 固定输入对象、来源、访问控制和权威水位快照。
\item 离线流程执行有界回放、比较、反例搜索和关联探索。
\item 结果首先形成假设。
\item 提交前重新核验来源、当前修订、权限和清除状态。
\item 目标领域所有者决定接受、拒绝、标记冲突或判定过期。
\end{enumerate}

\section{行为反馈}

动机目标、记忆、自我和社会认知可以作为 Heptalogos 的行为决策输入。Heptalogos 形成的承诺、执行结果和外部效果通过类型明确的结果回执返回 Nous Wave。回执本身形成来源事实或使用结果，再由认知领域决定是否形成长期修订。
